% ============================================================
\chapter{S\'eries Temporelles Multivari\'ees --- VAR et Co\"int\'egration}
\label{ch:var}
% ============================================================

\section{Exemple motivant}
Le taux d'int\'er\^et et l'inflation \'evoluent conjointement : la hausse de l'un influence l'autre. Un mod\`ele \textbf{VAR} capture ces interd\'ependances dynamiques. La \textbf{co\"int\'egration} r\'ev\`ele des relations d'\'equilibre long terme entre variables non stationnaires.

\section{Processus vectoriels stationnaires}

\begin{definition}[Processus vectoriel stationnaire]
$(\bm{Y}_t)_{t\in\Z}$ avec $\bm{Y}_t\in\R^k$ est \textbf{faiblement stationnaire} si :
\begin{enumerate}
\item $\E[\|\bm{Y}_t\|^2]<\infty$,
\item $\bm\mu = \E[\bm{Y}_t]$ ne d\'epend pas de $t$,
\item $\Gamma(h) = \Cov(\bm{Y}_t,\bm{Y}_{t+h}) = \E[(\bm{Y}_t-\bm\mu)(\bm{Y}_{t+h}-\bm\mu)^\top]$ ne d\'epend que de $h$.
\end{enumerate}
Note : $\Gamma(h)=\Gamma(-h)^\top$ (mais $\Gamma(h)\neq\Gamma(h)^\top$ en g\'en\'eral).
\end{definition}

\section{Mod\`ele VAR$(p)$}\index{VAR}

\begin{definition}[VAR$(p)$]
\[
\bm{Y}_t = \bm{c} + A_1\bm{Y}_{t-1} + A_2\bm{Y}_{t-2} + \cdots + A_p\bm{Y}_{t-p} + \bm{u}_t,
\]
o\`u $\bm{u}_t\sim\mathrm{BB}(\bm{0},\Sigma_u)$ ($k$-dimensionnel) et $A_i\in\R^{k\times k}$.
\end{definition}

\begin{theoreme}[Stationnarit\'e d'un VAR$(p)$]
Le VAR$(p)$ est stationnaire si et seulement si toutes les racines de
\[
\det(I_k - A_1 z - \cdots - A_p z^p) = 0
\]
sont en dehors du disque unit\'e.
\end{theoreme}

\subsection{Estimation}

\begin{theoreme}
Les param\`etres d'un VAR$(p)$ peuvent \^etre estim\'es \'equation par \'equation par OLS. L'estimateur est consistant et asymptotiquement normal.
\end{theoreme}

\section{Causalit\'e de Granger}\index{causalite de Granger@causalit\'e de Granger}

\begin{definition}[Causalit\'e de Granger]
$Y_{2,t}$ \textbf{cause au sens de Granger} $Y_{1,t}$ si les valeurs pass\'ees de $Y_2$ am\'eliorent la pr\'evision de $Y_1$, i.e.\ si dans le VAR, les coefficients associ\'es aux retards de $Y_2$ dans l'\'equation de $Y_1$ sont conjointement significatifs.

Test : $H_0 : A_1^{(12)}=\cdots=A_p^{(12)}=0$ par test de Wald ($F$-test).
\end{definition}

\section{Fonctions de r\'eponse impulsionnelle}\index{reponse impulsionnelle@r\'eponse impulsionnelle}

\begin{definition}[IRF]
La \textbf{fonction de r\'eponse impulsionnelle} mesure l'effet d'un choc unitaire sur $u_{j,t}$ sur $Y_{i,t+h}$ :
\[
\Psi_h = \frac{\partial\bm{Y}_{t+h}}{\partial\bm{u}_t^\top},
\]
o\`u $\bm{Y}_t = \sum_{h=0}^{\infty}\Psi_h\bm{u}_{t-h}$ est la repr\'esentation MA$(\infty)$.
\end{definition}

\section{Co\"int\'egration}\index{cointegration@co\"int\'egration}

\begin{definition}[Co\"int\'egration]
Des s\'eries $I(1)$ $Y_{1,t},\ldots,Y_{k,t}$ sont \textbf{co\"int\'egr\'ees} s'il existe $\bm\beta\neq\bm{0}$ tel que $\bm\beta^\top\bm{Y}_t\sim I(0)$. Le vecteur $\bm\beta$ est un \textbf{vecteur de co\"int\'egration} et $\bm\beta^\top\bm{Y}_t$ repr\'esente une relation d'\'equilibre long terme.
\end{definition}

\begin{exemple}
Les prix spot et futures d'une mati\`ere premi\`ere sont chacun $I(1)$, mais leur diff\'erence (la base) est $I(0)$ : ils sont co\"int\'egr\'es avec $\bm\beta=(1,-1)^\top$.
\end{exemple}

\begin{definition}[Mod\`ele \`a correction d'erreur (VECM)]\index{VECM}
\[
\nabla\bm{Y}_t = \Pi\bm{Y}_{t-1} + \sum_{i=1}^{p-1}\Gamma_i\nabla\bm{Y}_{t-i} + \bm{u}_t,
\]
o\`u $\Pi=\alpha\beta^\top$ avec $\alpha\in\R^{k\times r}$ (vitesses d'ajustement) et $\beta\in\R^{k\times r}$ (vecteurs de co\"int\'egration), $r=\rang(\Pi)$ est le nombre de relations de co\"int\'egration.
\end{definition}

\begin{theoreme}[Repr\'esentation de Granger]
Si $\bm{Y}_t\sim I(1)$ et co\"int\'egr\'e de rang $r$ :
\begin{itemize}
\item $0<r<k$ : co\"int\'egration. Le VECM est la bonne sp\'ecification.
\item $r=0$ : pas de co\"int\'egration, VAR en diff\'erences.
\item $r=k$ : toutes les s\'eries sont $I(0)$, VAR en niveaux.
\end{itemize}
\end{theoreme}

\subsection{Test de Johansen}\index{test!Johansen}

\begin{definition}
Le test de Johansen d\'etermine le rang de co\"int\'egration $r$ par des tests s\'equentiels de la valeur propre maximale ou de la trace.
\end{definition}

\section{Impl\'ementation}

\begin{codeblock}[title=\textbf{Python}]
\begin{minted}{python}
import numpy as np
import pandas as pd
import matplotlib.pyplot as plt
from statsmodels.tsa.api import VAR
from statsmodels.tsa.vector_ar.vecm import coint_johansen

# Simuler un VAR(1) bivarié
np.random.seed(42)
n = 500
A = np.array([[0.5, 0.3], [-0.2, 0.6]])
Y = np.zeros((n, 2))
for t in range(1, n):
    Y[t] = A @ Y[t-1] + np.random.multivariate_normal([0,0],
                                      [[1, 0.3],[0.3, 1]])

df = pd.DataFrame(Y, columns=['Y1', 'Y2'])

# Ajuster VAR
model = VAR(df)
results = model.fit(maxlags=5, ic='bic')
print(results.summary())

# Causalité de Granger
granger = results.test_causality('Y1', 'Y2', kind='f')
print(f"\nGranger Y2 -> Y1: p-value = {granger.pvalue:.4f}")

# IRF
irf = results.irf(20)
irf.plot(orth=True)
plt.savefig('ch09_irf.pdf')
plt.show()

# Coïntégration (Johansen) sur données simulées I(1)
rw1 = np.cumsum(np.random.normal(0, 1, n))
rw2 = 0.5 * rw1 + np.cumsum(np.random.normal(0, 0.3, n))
data_coint = np.column_stack([rw1, rw2])
joh = coint_johansen(data_coint, det_order=0, k_ar_diff=1)
print("\nJohansen trace stats:", joh.lr1)
print("Critical values (90%, 95%, 99%):", joh.cvt)
\end{minted}
\end{codeblock}

\section{Exercices}

\begin{exercice}[$\star$]
V\'erifier que le VAR(1) $\bm{Y}_t = A\bm{Y}_{t-1}+\bm{u}_t$ avec $A=\begin{pmatrix}0{,}5&0\\0&0{,}3\end{pmatrix}$ est stationnaire. Calculer $\Gamma(0)$.
\end{exercice}

\begin{exercice}[$\star\star$]
Montrer que dans un VAR(1), $\Gamma(h)=A^h\Gamma(0)$ pour $h\geq 0$.
\end{exercice}

\begin{exercice}[$\star\star$]
Montrer que si $Y_{1,t}$ et $Y_{2,t}$ sont des marches al\'eatoires ind\'ependantes, elles ne sont pas co\"int\'egr\'ees.
\end{exercice}

\begin{exercice}[$\star\star\star$ -- Projet]
Estimer un VAR pour le couple (taux d'int\'er\^et, inflation) d'un pays. Tester la causalit\'e de Granger dans les deux sens et interpr\'eter les fonctions de r\'eponse impulsionnelle.
\end{exercice}

\begin{formules}
\begin{align*}
\text{VAR}(p) &: \bm{Y}_t = \bm{c}+A_1\bm{Y}_{t-1}+\cdots+A_p\bm{Y}_{t-p}+\bm{u}_t \\
\text{VECM} &: \nabla\bm{Y}_t=\alpha\beta^\top\bm{Y}_{t-1}+\sum\Gamma_i\nabla\bm{Y}_{t-i}+\bm{u}_t \\
\text{Co\"int\'egration} &: \bm\beta^\top\bm{Y}_t\sim I(0)\text{ avec }\bm{Y}_t\sim I(1)
\end{align*}
\end{formules}
